% Part: proof-theory
% Chapter: sequent-calculus
% Section: rules-proofs

\documentclass[../../../include/open-logic-section]{subfiles}

\begin{document}

\olfileid[hi]{pt}{seq}{rp}

\olsection{नियम और \usetoken{P}{proof}}

हम कुछ परिभाषाएँ देकर और कुछ पारिभाषिक शब्द निश्चित करके आरंभ करते हैं।

उदाहरण के लिए \Log{G3c} के उन दो नियमों पर विचार करें जो $!A$ और~$!B$ वाले
अनुक्रमों से $!A \land !B$ वाले अनुक्रमों को !!{proving} करने देते हैं:
\[
\Axiom$ !A, !B, \Gamma \fCenter \Delta$
\RightLabel{\LeftR{\land}}
\UnaryInf$ !A \land !B, \Gamma \fCenter \Delta$
\DisplayProof
\qquad
\Axiom$\Gamma \fCenter \Delta, !A$
\Axiom$ \Gamma \fCenter \Delta, !B$
\RightLabel{\RightR{\land}}
\BinaryInf$ \Gamma \fCenter \Delta, !A \land !B $
\DisplayProof
\]
रेखा के ऊपर के अनुक्रमों को \emph{आधार-वाक्य} और रेखा के नीचे के अनुक्रम को
\emph{निष्कर्ष} कहते हैं। प्रत्येक नियम में $\Gamma$ और $\Delta$ के
!!{formula}s को \emph{प्रसंग}, अथवा \emph{पार्श्व !!{formula}s}, कहते हैं।
निष्कर्ष में जो !!{formula} प्रसंग का भाग नहीं है उसे \emph{प्रमुख}
!!{formula}, और आधार-वाक्यों में प्रसंग के बाहर के !!{formula}s को
\emph{सक्रिय} (अथवा \emph{सहायक}) !!{formula}s कहते हैं।

परिमाणक नियम कुछ अधिक जटिल हैं। उदाहरणतः \Log{G1c} में $\lforall$ के नियम हैं:
\[\Axiom$ !A(t), \Gamma \fCenter \Delta$
\RightLabel{\LeftR{\lforall}}
\UnaryInf$ \lforall[x][!A(x)],\Gamma \fCenter \Delta$
\DisplayProof
\qquad
\Axiom$ \Gamma \fCenter \Delta, !A(a) $
\RightLabel{\RightR{\lforall}}
\UnaryInf$ \Gamma \fCenter \Delta, \lforall[x][!A(x)]$
\DisplayProof
\]
\LeftR{\lforall} नियम में सक्रिय !!{formula}~$!A(t)$ है, जहाँ $t$ कोई बंद
पद, अर्थात चरों से रहित पद, है। ऐसा निगमन सही है---अर्थात वह योजनाबद्ध नियम
\LeftR{\lforall} से मेल खाता है---यदि कोई !!{formula}~$!A$, कोई चर~$x$ और
कोई बंद पद~$t$ ऐसे हैं कि निष्कर्ष का प्रमुख !!{formula}
$\lforall[x][!A]$ और आधार-वाक्य का सक्रिय !!{formula}~$\Subst{!A}{t}{x}$
है। \RightR{\lforall} नियम इसी प्रकार काम करता है, बस किसी भी पद~$t$ के
बदले आधार-वाक्य का सक्रिय !!{formula}~$\Subst{!A}{a}{x}$ होना चाहिए, जहाँ
$a$ कोई !!a{constant} है। इस अचर को \RightR{\lforall} निगमन का
\emph{आइगेनचर} कहते हैं। निगमन तभी सही है जब नीचे वाले अनुक्रम में~$a$ न आए,
अर्थात $a$, $\Gamma$, $\Delta$ या~$!A$ में न आए। इसे \emph{आइगेनचर प्रतिबंध}
कहते हैं।

\Log{G1c} और \Log{G3c} जैसी अनुक्रम प्रणालियों में ऐसे योजनाबद्ध रूप से
निर्दिष्ट नियमों का संग्रह और कुछ (वे भी योजनाबद्ध रूप से निर्दिष्ट) अनुक्रम
होते हैं जिन्हें अभिगृहीतियाँ माना जाता है। ऐसी प्रणाली में !!^{proof}s,
अनुक्रमों के सीमित वृक्ष हैं जो निगमन नियमों का उपयोग करके अभिगृहीतियों से
आगमनतः उत्पन्न होते हैं। प्रत्येक पत्ता अभिगृहीति है और प्रत्येक अन्य अनुक्रम
उसके ठीक ऊपर के अनुक्रमों से प्रणाली के किसी निगमन नियम द्वारा मिलता है।

\begin{defn}\ollabel{defn:proof}
किसी अनुक्रम प्रणाली में \emph{!!^{proof}s}, निम्न प्रकार आगमनतः परिभाषित
अनुक्रमों के सीमित वृक्ष हैं:
\begin{enumerate}
  \item\ollabel{defn:prf:axiom} प्रत्येक अभिगृहीति !!a{proof} है।
  \item\ollabel{defn:prf:one-prem} यदि $\pi_1$, अनुक्रम~$S_1$ पर समाप्त
  होने वाला !!a{proof} है और $S$, एक आधार-वाक्य वाले नियम~$R$ द्वारा~$S_1$
  से मिलने वाला अनुक्रम है, तो
  \[
  \AxiomC{}
  \RightLabel{$\pi_1$}
  \DeduceC{$S_1$}
  \RightLabel{$R$}
  \UnaryInfC{$S$}
  \DisplayProof
  \]
  !!a{proof} है।
  \item\ollabel{defn:prf:two-prem} यदि $\pi_1$ और $\pi_2$, क्रमशः अनुक्रमों
  ~$S_1$ और~$S_2$ पर समाप्त होने वाले !!{proof}s हैं और $S$, दो आधार-वाक्यों
  वाले नियम~$R$ द्वारा~$S_1$ और $S_2$ से मिलने वाला अनुक्रम है, तो
  \[
  \AxiomC{}
  \RightLabel{$\pi_1$}
  \DeduceC{$S_1$}
  \AxiomC{}
  \RightLabel{$\pi_2$}
  \DeduceC{$S_2$}
  \RightLabel{$R$}
  \BinaryInfC{$S$}
  \DisplayProof
  \]
  !!a{proof} है।
\end{enumerate}
\end{defn}

\Olref{defn:proof}, !!{proof}s को अनुक्रमों के वृक्षों के रूप में परिभाषित
करता है। हम इन वृक्षों को ``ऊपर की ओर'' बढ़ता हुआ समझते हैं। वृक्ष के सबसे
ऊपरी (पत्ता) नोड अभिगृहीतियाँ हैं और उन्हें \emph{आरंभिक अनुक्रम} कहते हैं।
सबसे नीचे वाले अनुक्रम को \emph{अंतिम अनुक्रम} कहते हैं। यदि $\pi$, अंतिम
अनुक्रम~$S$ वाला !!a{proof} है, तो हम $\pi \Proves S$ भी लिखते हैं। यदि
किसी अनुक्रम~$S$ का !!a{proof} है, तो हम $\Proves S$ लिखते हैं; कभी-कभी
$\Proves$ चिह्न की बाईं ओर वह प्रणाली भी दिखाते हैं जिसमें प्रमाण है, जैसे
$\Log{G1c} \Proves !A, !B \Sequent !A \land !B$। आरंभिक अनुक्रमों को छोड़कर
!!a{proof} का प्रत्येक अनुक्रम किसी नियम~$R$ द्वारा किए गए निगमन का
\emph{निष्कर्ष} है। !!a{proof} में किसी अनुक्रम के ठीक ऊपर के एक या दो
अनुक्रमों (अर्थात उस नोड के जनकों) को निगमन के \emph{आधार-वाक्य} कहते हैं।

!!a{proof} की आगमनात्मक रचना में उपयोग हुए !!{proof}s को उसके
\emph{उप-!!{proof}s} कहते हैं। किसी निगमन के आधार-वाक्यों पर समाप्त होने वाले
उप-!!{proof}s को उस निगमन के \emph{निकटस्थ उप-!!{proof}s} कहते हैं।

प्रायः !!{proof}s की जटिलता को किसी प्राकृतिक संख्या से मापना आवश्यक होता
है। हम !!a{proof} में अनुक्रमों की संख्या मात्र गिन सकते हैं, किंतु अक्सर
!!a{proof} की !!{height} अधिक उपयोगी माप है: यह अंतिम अनुक्रम और किसी आरंभिक
अनुक्रम के बीच अनुक्रमों की अधिकतम संख्या है। हम इसकी आगमनात्मक परिभाषा भी
दे सकते हैं।

\begin{defn}
!!a{proof}~$\pi$ की \emph{!!{height}}~$\pheight{\pi}$ इस प्रकार आगमनतः
परिभाषित है।
\begin{enumerate}
  \item यदि $\pi$ में केवल एक अभिगृहीति है, तो $\pheight{\pi} = 0$।
  \item यदि $\pi$ एक आधार-वाक्य वाले ऐसे निगमन पर समाप्त होता है जिसका
  निकटस्थ उप-!!{proof}~$\pi_1$ है, तो $\pheight{\pi} = \pheight{\pi_1} + 1$।
  \item यदि $\pi$ दो आधार-वाक्यों वाले ऐसे निगमन पर समाप्त होता है जिसके
  निकटस्थ उप-!!{proof}s~$\pi_1$ और~$\pi_2$ हैं, तो $\pheight{\pi} =
  \max(\pheight{\pi_1}, \pheight{\pi_2}) + 1$.
\end{enumerate}
यदि अनुक्रम~$S$ का ऊँचाई $\le n$ वाला !!a{proof} है, तो हम $\Proves[n] S$
लिखते हैं।
\end{defn}

\begin{ex}
\Log{G3c} में $!C, \Gamma \Sequent \Delta, !C$ रूप का हर अनुक्रम अभिगृहीति
है, जहाँ $!C$ परमाण्विक होना चाहिए। मान लें कि $!D$ और $!E$ परमाण्विक वाक्य
हैं। तब $!D, !E \Sequent !D$ और $!D, !E \Sequent !E$, अभिगृहीति $!C,
\Gamma \Sequent \Delta, !C$ के दृष्टांत हैं। उदाहरणतः यदि $!C, \Gamma
\Sequent \Delta, !C$ में $!C \ident !E$, $\Gamma = \{!D\}$ और $\Delta =
\emptyset$ लें, तो ठीक अनुक्रम $!D, !E \Sequent !E$ मिलता है। (याद रखें कि
हम बहु\emph{समुच्चयों} से काम कर रहे हैं, इसलिए $!D, !E \Sequent !E$ और
$!E, !D \Sequent !D$ एक ही अनुक्रम हैं।)

चूँकि $!D, !E \Sequent !D$, \Log{G3c} की किसी अभिगृहीति का दृष्टांत है, वह
अपने आप में \Log{G3c} का !!a{proof} माना जाता है; और
\olref{defn:proof}\olref{defn:prf:axiom} से $!D, !E \Sequent !E$ भी। इन्हें
$\pi_1$ और~$\pi_2$ कहें। \olref{defn:prf:two-prem} से हम इन दोनों !!{proof}s
को लेकर दो आधार-वाक्यों वाले नियम~$\RightR{\land}$ से नया प्रमाण ($\pi_3$)
बना सकते हैं:
\[
\Axiom$!D, !E \fCenter !D$
\Axiom$!D, !E \fCenter !E$
\RightLabel{\RightR{\land}}
\BinaryInf$!D, !E  \fCenter !D \land !E $
\DisplayProof
\]
$\pi_3$ से फिर हमें यह !!{proof}~$\pi_4$ मिलता है:
\[
\Axiom$!D, !E \fCenter !D$
\Axiom$!D, !E \fCenter !E$
\RightLabel{\RightR{\land}}
\insertBetweenHyps{\hspace{5ex}}
\BinaryInf$!D, !E  \fCenter !D \land !E $
\LeftSubproofLabel{$\pi_3$}
\RightLabel{\LeftR{\land}}
\UnaryInf$!D \land !E  \fCenter !D \land !E$
\RightSubproofLabel{$\pi_4$}
\DisplayProof
\]
इसके लिए एक आधार-वाक्य वाले नियम~$\LeftR{\land}$ का उपयोग हुआ है
(\olref{defn:proof}\olref{defn:prf:one-prem} का उपयोग करके)। $\pi_4$ के
अंतिम निगमन का निकटस्थ उप-!!{proof}~$\pi_3$ है। $\RightR{\land}$ निगमन के
निकटस्थ उप-!!{proof}s $\pi_1$ और~$\pi_2$ हैं। $\pi_4$ के उप-!!{proof}s में
$\pi_1$, $\pi_2$, $\pi_3$ और स्वयं $\pi_4$, सभी हैं। हमारे पास
$\pheight{\pi_1} = \pheight{\pi_2} = 0$, $\pheight{\pi_3} = 1$ और
$\pheight{\pi_4} = 2$ हैं। हमने दिखाया है कि किन्हीं भी परमाण्विक $!D$ और
~$!E$ के लिए $\Log{G3c} \Proves[4] !D \land !E \fCenter !D \land !E$।
\end{ex}

\end{document}
